\newpage

\begin{center}
  \textsc{Zusammenfassung}
\end{center}

\noindent

Kommunikation ist entscheidend im Wettbewerb um Ressourcen bei
Tieren, formt Dominanzhierarchien und kann die Notwendigkeit physischer
Auseinandersetzungen reduzieren. Die Funktion von \emph{chirps} – einer verbreiteten
Art elektrokommunikativer Signale bei elektrischen Fischen – ist 
umstritten. Bisherige Studien konzentrierten sich aufgrund der Schwierigkeiten
bei der Detektion von \emph{chirps} in freien Interaktionen vor allem auf isolierte
Fische. Diese Arbeit präsentiert eine \emph{Deep-Learning Pipeline} zur
automatisierten Erkennung von \emph{chirps} bei frei interagierenden
\textit{Apteronotus leptorhynchus}, um die Rolle von \emph{chirps} im ritualisierten
Wettbewerb um Unterschlupf und darüber hinaus zu erforschen.

Die entwickelte Pipeline nutzt \emph{Convolutional Neural Networks} für die
\emph{chirp}-Erkennung. Diese werden mit Spektrogrammen trainiert, auf denen
die \emph{chirps} markiert sind. Nach der Erkennung kommen speziell entwickelte
Algorithmen zum Einsatz, um die detektierten \emph{chirps} dem jeweils
kommunizierenden Fisch zuzuordnen. Diese Methode ermöglicht eine präzise
Identifizierung und Zuordnung von \emph{chirps} zu einzelnen Fischen und führte
zu einem Datensatz von über 73.000 \emph{chirps}, die bei freien Interaktionen
zwischen zwei Fischen produziert wurden.

In ersten Analysen wurden die Muster der \emph{chirps} und deren
Verhaltenskorrelationen untersucht. Eine wichtige Erkenntnis war die starke
Korrelation zwischen \emph{chirps} von submissiven Individuen und dem Abbruch
aggressiver Interaktionen bei Wettbewerben um Unterschlupf. Diese Beobachtung
legt nahe, dass \emph{chirps} über die sensorische Wahrnehmung hinaus eine
bedeutende Rolle spielen. Die von submissiven Fischen ausgesendeten
\emph{chirps} schienen jedoch keinen messbaren Einfluss auf das Verhalten der
dominanten Fische zu haben, was die Notwendigkeit weiterer Untersuchungen
hervorhebt.

Darüber hinaus lässt sich diese automatisierte Methode auch auf Felddaten
anwenden. Vorherige Studien haben gezeigt, dass die markerlose Verfolgung
mehrerer Individuen in ihrer natürlichen Umgebung über die Erfassung
elektrischer Felder auf Elektroden-Gittern möglich ist. Die \emph{Chirp
Detection Pipeline} baut auf diesen Erkenntnissen auf und eröffnet die
Möglichkeit, das gesamte elektrokommunikative Repertoire aus solchen Aufnahmen
zu extrahieren. Diese Entwicklung bietet erstmals die Gelegenheit, umfassende
Informationen über Identität, Geschlecht, Bewegungen und nun auch
Kommunikationssignale für quantitative ethologische Studien von ungestörten
Populationen anhand von Daten aus einer einzigen Modalität zu analysieren.

\newpage

\begin{center}
  \textsc{Abstract}
\end{center}

\noindent

Effective communication plays a crucial role in resource competition among
animals, shaping dominance hierarchies and reducing the necessity for physical
confrontations. The function of chirps, a common form of
electrocommunication signal in electric fish, is still a topic of debate. Past
studies have mainly focused on isolated fish due to the technical limitations
in detecting chirps during interactions. This thesis introduces a deep learning
pipeline for automated chirp detection in freely moving Brown Ghost Knifefish
to investigate the function of chirps in staged competition scenarios and
beyond.

The pipeline applies deep convolutional neural networks for chirp detection,
which are fine-tuned using labeled spectrograms of chirping fish. Following
detection, it uses custom algorithms to link detected chirps with the specific
fish producing them. This methodology facilitated the identification and
allocation of chirps to individual fish, resulting in a substantial dataset of
over 73,000 chirps detected during unconstrained interactions.

Initial analyses focused on the patterns of chirping and their behavioral
associations. A notable finding is a significant correlation between chirps
emitted by subordinate individuals and the cessation of aggressive interactions
when competing for shelter. This observation strongly implies that chirps serve
a purpose beyond active sensing. Still, the chirps emitted by the subordinate
fish did not seem to influence the behavior of the dominant fish in a
measurable manner, which highlights the need for further investigation.

Beyond laboratory settings, this automated methodology can be applied to field
recordings as well. Prior research has illustrated the feasibility of
markerless tracking of multiple individuals in their natural environments by
capturing their electric fields on electrode grids. Building on this
foundation, this approach provides the capability to extract the complete
electrocommunication repertoire from such recordings. This advancement now
represents the opportunity to extract information concerning identity, sex,
movements, and communication signals for quantitative ethological studies in
undisturbed populations using data from a single modality for the first time.
